\documentclass[11pt,a4paper]{article}

% --------------------------------------------------
% PAQUETES
% --------------------------------------------------
\usepackage[utf8]{inputenc}
\usepackage[T1]{fontenc}
\IfFileExists{spanish.ldf}{%
  \usepackage[spanish,es-tabla]{babel}%
}{%
  \usepackage[english]{babel}%
}
\usepackage{lmodern}
\usepackage{geometry}
\usepackage{graphicx}
\usepackage{fancyhdr}
\usepackage{xcolor}
\usepackage{titlesec}
\usepackage[nopatch=footnote]{microtype}
\usepackage{hyperref}
\usepackage{enumitem}
\usepackage{array}
\usepackage{tabularx}
\usepackage{booktabs}
\usepackage{longtable}
\usepackage{ragged2e}
\usepackage{float}
\usepackage{caption}
\usepackage{setspace}
\usepackage{colortbl}
\usepackage{lastpage}
\usepackage{pdflscape}
\usepackage{pifont}
\usepackage{url}

% --------------------------------------------------
% CONFIGURACION GENERAL
% --------------------------------------------------
\geometry{
  top=2.6cm,
  bottom=2.6cm,
  left=2.6cm,
  right=2.6cm,
  headheight=15pt,
  footskip=1.2cm
}

\definecolor{azulprincipal}{HTML}{1F3C88}
\definecolor{azulsuave}{HTML}{EAF0FB}
\definecolor{grislinea}{HTML}{D9DEE8}
\definecolor{gristexto}{HTML}{333333}
\definecolor{rojopendiente}{HTML}{B00020}
\definecolor{verdesuave}{HTML}{EAF7EA}
\definecolor{amarillosuave}{HTML}{FFF7E0}

\hypersetup{
  colorlinks=true,
  linkcolor=azulprincipal,
  urlcolor=azulprincipal,
  citecolor=azulprincipal,
  pdftitle={Reporte final de sprints, milestones y postmortem},
  pdfauthor={SYNTIX TECH - Pontificia Universidad Javeriana}
}

\setlength{\parindent}{0pt}
\setlength{\parskip}{0.6em}
\setstretch{1.12}
\renewcommand{\arraystretch}{1.25}
\setlength{\emergencystretch}{2em}

\setcounter{secnumdepth}{4}
\setcounter{tocdepth}{3}
\renewcommand{\contentsname}{Tabla de contenido}
\AtBeginDocument{%
  \renewcommand{\contentsname}{Tabla de contenido}%
  \renewcommand{\tablename}{Tabla}%
  \renewcommand{\figurename}{Figura}%
}

% --------------------------------------------------
% FORMATO DE TITULOS
% --------------------------------------------------
\titleformat{\section}
  {\normalfont\Large\bfseries\color{azulprincipal}}
  {\thesection.}{0.6em}{}

\titleformat{\subsection}
  {\normalfont\large\bfseries\color{azulprincipal}}
  {\thesubsection.}{0.6em}{}

\titleformat{\subsubsection}
  {\normalfont\normalsize\bfseries\color{azulprincipal}}
  {\thesubsubsection.}{0.6em}{}

\titleformat{\paragraph}[block]
  {\normalfont\normalsize\bfseries\color{azulprincipal}}
  {\theparagraph.}{0.6em}{}

\titlespacing*{\section}{0pt}{1.6em}{0.8em}
\titlespacing*{\subsection}{0pt}{1.2em}{0.6em}
\titlespacing*{\subsubsection}{0pt}{1em}{0.5em}
\titlespacing*{\paragraph}{0pt}{0.9em}{0.4em}

% --------------------------------------------------
% ENCABEZADO Y PIE DE PAGINA
% --------------------------------------------------
\pagestyle{fancy}
\fancyhf{}
\fancyhead[L]{\textcolor{azulprincipal}{Reporte final Agile/Postmortem}}
\fancyhead[R]{\textcolor{azulprincipal}{DriveControl / AutoMinder Enterprise}}
\fancyfoot[C]{\textcolor{gristexto}{Pagina \thepage\ de \pageref*{LastPage}}}
\renewcommand{\headrulewidth}{0.6pt}
\renewcommand{\footrulewidth}{0pt}

\captionsetup{
  font=small,
  labelfont=bf
}

\newcolumntype{Y}{>{\RaggedRight\arraybackslash}X}
\newcolumntype{P}[1]{>{\RaggedRight\arraybackslash}p{#1}}

\newcommand{\code}[1]{\texttt{\detokenize{#1}}}

\begin{document}

% --------------------------------------------------
% PORTADA
% --------------------------------------------------
\hypersetup{pageanchor=false}
\begin{titlepage}
\centering
\vspace*{1.2cm}
{\Large \textbf{Pontificia Universidad Javeriana}}\\[0.4em]
{\large Departamento de Ingenieria de Sistemas}\\[2.4cm]

{\Huge \textbf{Reporte final de sprints, milestones y postmortem}}\\[0.7em]
{\Large \textbf{DriveControl / AutoMinder Enterprise}}\\[1.2cm]

\begin{tabular}{rl}
\textbf{Proyecto:} & DriveControl / AutoMinder Enterprise \\
\textbf{Equipo:} & SYNTIX TECH \\
\textbf{Curso:} & Fundamentos de Ingenieria de Software \\
\textbf{Profesor:} & Luis Gabriel Moreno Sandoval \\
\textbf{Fecha de consolidacion:} & 19 de mayo de 2026 \\
\textbf{Repositorio:} & \url{https://github.com/puj-course/FIS_2610_3517_G4} \\
\end{tabular}

\vfill

\begin{tabular}{lll}
\toprule
\textbf{Nombre} & \textbf{Usuario} & \textbf{Rol} \\
\midrule
Sebastian Ramirez Maldonado & @Sarm-m & Scrum Master \\
Samuel Freile & @samuelfl680 & Configuration Manager \\
Sebastian Rodriguez Ramirez & @Juserora & Quality Assurance Lead \\
Solon Losada & @solonlosada2006 & DevOps Engineer \\
Sebastian Vargas & @juanvargax & Product Owner / Sprint Planner \\
\bottomrule
\end{tabular}

\vfill
{\small Informe generado a partir del reporte general en Markdown y adaptado al formato LaTeX academico.}
\end{titlepage}
\hypersetup{pageanchor=true}

\pagenumbering{roman}
\tableofcontents
\clearpage

\pagenumbering{arabic}

\section*{Introduccion}
\addcontentsline{toc}{section}{Introduccion}

Este informe general consolida el cierre metodologico y tecnico del semestre para el proyecto \textbf{DriveControl / AutoMinder Enterprise}. Su proposito es reunir en un solo documento la evidencia de Scrum academico, milestones, sprints, issues, historias de usuario, Pull Requests, commits, revisiones, control de versiones y postmortem general del equipo SYNTIX TECH.

El documento conserva el contenido del reporte general en Markdown y lo organiza en formato LaTeX para facilitar su entrega academica, su incorporacion en Overleaf y su lectura como informe final. La estructura mantiene las cifras verificadas, las aclaraciones sobre datos pendientes de reconciliacion y la relacion explicita con la rubrica de metodologias agiles y postmortem.

\section{1. Resumen ejecutivo}

Este documento es la fuente central de la evidencia ágil del semestre. Consolida Scrum académico, GitHub Issues, milestones, Pull Requests, ramas, commits, revisiones disponibles, trazabilidad y postmortem técnico/de proceso para cubrir explícitamente el criterio \textbf{Metodologías ágiles y postmortem --- 30\%}.

La fuente principal para las métricas nuevas es GitHub live mediante \texttt{\detokenize{gh}}, complementada por Git local y anexos versionados en \texttt{\detokenize{docs/Entrega-Final/anexos/}}. Las cifras antiguas del reporte previo se conservan como histórico cuando no coinciden con GitHub live; no se usan para inflar artificialmente el desempeño.

\section{2. Fuentes de datos usadas}

\begingroup
\small
\setlength{\LTpre}{0.4em}
\setlength{\LTpost}{0.8em}
\begin{longtable}{>{\RaggedRight\arraybackslash}p{0.313\linewidth}>{\RaggedRight\arraybackslash}p{0.313\linewidth}>{\RaggedRight\arraybackslash}p{0.313\linewidth}}
\toprule
\textbf{Fuente} & \textbf{Uso} & \textbf{Estado} \\
\midrule
\endfirsthead
\toprule
\textbf{Fuente} & \textbf{Uso} & \textbf{Estado} \\
\midrule
\endhead
\midrule
\multicolumn{3}{r}{\small\textit{Continua en la siguiente pagina}} \\
\endfoot
\bottomrule
\endlastfoot
GitHub CLI autenticado sobre \texttt{\detokenize{puj-course/FIS_2610_3517_G4}} & Issues, PRs, reviews, milestones & Verificado \\
Git local \texttt{\detokenize{git log --all}} & Commits por sprint y por integrante & Verificado \\
\texttt{\detokenize{docs/Entrega-Final/anexos/issues.json}} & Respaldo de issues exportadas & Disponible \\
\texttt{\detokenize{docs/Entrega-Final/anexos/pull_requests.json}} & Respaldo de PRs exportados & Disponible \\
\texttt{\detokenize{docs/Entrega-Final/anexos/commits_detalle.tsv}} & Respaldo de commits exportados & Disponible \\
\texttt{\detokenize{docs/Entrega-Final/evidencias/}} & Capturas de milestones, tablero, contributors, PRs y commits & Disponible parcialmente \\
\texttt{\detokenize{Sesión 2- Postmortem.pdf}} & Guía del profesor sobre producto, esfuerzo, proceso, datos de calidad, roles y Starfish & Verificado \\
\end{longtable}
\endgroup

\section{3. Milestones válidos del semestre}

Solo se usan estos cuatro milestones. Se ignora explícitamente el milestone de taller \texttt{\detokenize{Taller, para aprender github}}.

\begin{landscape}
\begingroup
\scriptsize
\setlength{\LTpre}{0.4em}
\setlength{\LTpost}{0.8em}
\begin{longtable}{>{\RaggedRight\arraybackslash}p{0.153\linewidth}>{\RaggedRight\arraybackslash}p{0.153\linewidth}>{\RaggedRight\arraybackslash}p{0.153\linewidth}>{\RaggedRight\arraybackslash}p{0.153\linewidth}>{\RaggedRight\arraybackslash}p{0.153\linewidth}>{\RaggedRight\arraybackslash}p{0.153\linewidth}}
\toprule
\textbf{Milestone} & \textbf{Objetivo} & \textbf{Issues cerradas} & \textbf{HU cerradas trazables} & \textbf{PRs con milestone GitHub} & \textbf{Estado} \\
\midrule
\endfirsthead
\toprule
\textbf{Milestone} & \textbf{Objetivo} & \textbf{Issues cerradas} & \textbf{HU cerradas trazables} & \textbf{PRs con milestone GitHub} & \textbf{Estado} \\
\midrule
\endhead
\midrule
\multicolumn{6}{r}{\small\textit{Continua en la siguiente pagina}} \\
\endfoot
\bottomrule
\endlastfoot
Milestone 1 - Base funcional del sistema & Establecer base MERN, autenticación, dashboard inicial y GitFlow. & 76 & 24 & 9 & Finalizado \\
Milestone 2 - Gestión básica de flota & Consolidar gestión de flota, patrones y validación RUNT. & 37 & 14 & 1 & Finalizado \\
Milestone 3 - Gestión documental y monitoreo & Integrar documentos, persistencia, monitoreo y controles de seguridad. & 38 & 14 & 0 & Finalizado \\
Milestone 4 - Dashboard, alertas y cierre del sistema & Cerrar dashboard, alertas, CI/CD, SMS, calidad y entrega final. & 132 & 20 & 3 & Finalizado con PR abierto de cierre \\
\end{longtable}
\endgroup
\end{landscape}

Lectura para la rúbrica: los 4 milestones funcionales aparecen cerrados en GitHub para issues. El único elemento abierto asociado a M4 es el PR \texttt{\detokenize{#628}}, por lo que se reporta como pendiente de integración y no como issue funcional abierta.

\section{4. Métricas consolidadas}

\begingroup
\small
\setlength{\LTpre}{0.4em}
\setlength{\LTpost}{0.8em}
\begin{longtable}{>{\RaggedRight\arraybackslash}p{0.313\linewidth}>{\RaggedRight\arraybackslash}p{0.313\linewidth}>{\RaggedRight\arraybackslash}p{0.313\linewidth}}
\toprule
\textbf{Métrica} & \textbf{Valor verificado} & \textbf{Interpretación} \\
\midrule
\endfirsthead
\toprule
\textbf{Métrica} & \textbf{Valor verificado} & \textbf{Interpretación} \\
\midrule
\endhead
\midrule
\multicolumn{3}{r}{\small\textit{Continua en la siguiente pagina}} \\
\endfoot
\bottomrule
\endlastfoot
Sprints académicos revisados & 13 & Cobertura del semestre completo, no solo Sprint 13. \\
Issues cerradas en M1-M4 & 283 & Issues con milestone funcional válido y estado cerrado. \\
HU cerradas trazables & 72 & Conteo estricto por patrón \texttt{\detokenize{HU-}}, \texttt{\detokenize{HU }}, \texttt{\detokenize{[HU-...]}} o \texttt{\detokenize{Historia de usuario}}. \\
Promedio HU/sprint & 5.54 & 72 HU / 13 sprints. \\
Commits en corte académico & 630 & \texttt{\detokenize{git log --all}} entre 2026-02-16 y 2026-05-17 23:59:59. \\
Promedio commits/sprint & 48.46 & 630 commits / 13 sprints. \\
Pull Requests totales & 161 & PRs en todos los estados. \\
Pull Requests mergeadas & 134 & Evidencia de integración continua por PR. \\
Pull Requests abiertas & 1 & PR \texttt{\detokenize{#628}}, asociado a M4. \\
Pull Requests cerradas sin merge & 26 & Trabajo descartado, reemplazado o no integrado. \\
Reviews detectadas & 8 & Dato disponible desde \texttt{\detokenize{gh pr list --json reviews}}. \\
Tags/releases verificables & 0 & No existen releases/tags verificables en GitHub local/live. \\
\end{longtable}
\endgroup

\section{5. Cifras históricas pendientes de reconciliación}

El reporte anterior contenía 88 HU, 375 issues y 405 commits. Esas cifras se conservan como histórico del documento previo, pero quedan marcadas como \textbf{pendientes de reconciliación} porque el cálculo live con GitHub y Git local arroja 72 HU estrictas, 283 issues cerradas en M1-M4 y 630 commits en el corte académico. La diferencia puede explicarse por criterios distintos: conteo manual, issues sin milestone, issues de taller, commits filtrados por rama o exclusión de merges.

\section{6. Trazabilidad por sprint}

\begin{landscape}
\begingroup
\scriptsize
\setlength{\LTpre}{0.4em}
\setlength{\LTpost}{0.8em}
\begin{longtable}{>{\RaggedRight\arraybackslash}p{0.153\linewidth}>{\RaggedRight\arraybackslash}p{0.153\linewidth}>{\RaggedRight\arraybackslash}p{0.153\linewidth}>{\RaggedRight\arraybackslash}p{0.153\linewidth}>{\RaggedRight\arraybackslash}p{0.153\linewidth}>{\RaggedRight\arraybackslash}p{0.153\linewidth}}
\toprule
\textbf{Sprint} & \textbf{Periodo} & \textbf{Issues cerradas M1-M4} & \textbf{HU cerradas estrictas} & \textbf{Commits Git local} & \textbf{Lectura} \\
\midrule
\endfirsthead
\toprule
\textbf{Sprint} & \textbf{Periodo} & \textbf{Issues cerradas M1-M4} & \textbf{HU cerradas estrictas} & \textbf{Commits Git local} & \textbf{Lectura} \\
\midrule
\endhead
\midrule
\multicolumn{6}{r}{\small\textit{Continua en la siguiente pagina}} \\
\endfoot
\bottomrule
\endlastfoot
S1 & 2026-02-16 a 2026-02-22 & 0 & 0 & 30 & Inicio técnico y preparación de base. \\
S2 & 2026-02-23 a 2026-03-01 & 4 & 0 & 14 & Primer cierre verificable de issues. \\
S3 & 2026-03-02 a 2026-03-08 & 16 & 7 & 33 & Consolidación inicial de HU. \\
S4 & 2026-03-09 a 2026-03-15 & 25 & 7 & 30 & Cierre funcional de M1. \\
S5 & 2026-03-16 a 2026-03-22 & 11 & 2 & 30 & Arranque de gestión de flota. \\
S6 & 2026-03-23 a 2026-03-29 & 10 & 2 & 26 & Avance incremental de flota. \\
S7 & 2026-03-30 a 2026-04-05 & 16 & 6 & 14 & Integración de validaciones. \\
S8 & 2026-04-06 a 2026-04-12 & 36 & 10 & 94 & Cierre fuerte de M2 y deuda técnica. \\
S9 & 2026-04-13 a 2026-04-19 & 42 & 14 & 76 & Mayor cierre funcional de M3. \\
S10 & 2026-04-20 a 2026-04-26 & 9 & 2 & 6 & Cierre documental y monitoreo. \\
S11 & 2026-04-27 a 2026-05-03 & 40 & 5 & 81 & Inicio del cierre técnico. \\
S12 & 2026-05-04 a 2026-05-10 & 2 & 2 & 61 & Preparación CI/CD, QA y evidencia. \\
S13 & 2026-05-11 a 2026-05-17 & 72 & 15 & 135 & Cierre de rúbrica, Docker, SonarCloud, SMS y documentación. \\
\end{longtable}
\endgroup
\end{landscape}

\section{7. Trazabilidad por milestone}

\begin{landscape}
\begingroup
\scriptsize
\setlength{\LTpre}{0.4em}
\setlength{\LTpost}{0.8em}
\begin{longtable}{>{\RaggedRight\arraybackslash}p{0.153\linewidth}>{\RaggedRight\arraybackslash}p{0.153\linewidth}>{\RaggedRight\arraybackslash}p{0.153\linewidth}>{\RaggedRight\arraybackslash}p{0.153\linewidth}>{\RaggedRight\arraybackslash}p{0.153\linewidth}>{\RaggedRight\arraybackslash}p{0.153\linewidth}}
\toprule
\textbf{Milestone} & \textbf{Sprints asociados} & \textbf{Commits por ventana de sprint} & \textbf{Issues cerradas} & \textbf{HU cerradas} & \textbf{PRs destacados} \\
\midrule
\endfirsthead
\toprule
\textbf{Milestone} & \textbf{Sprints asociados} & \textbf{Commits por ventana de sprint} & \textbf{Issues cerradas} & \textbf{HU cerradas} & \textbf{PRs destacados} \\
\midrule
\endhead
\midrule
\multicolumn{6}{r}{\small\textit{Continua en la siguiente pagina}} \\
\endfoot
\bottomrule
\endlastfoot
M1 & S1-S4 & 107 & 76 & 24 & \texttt{\detokenize{#130}}, \texttt{\detokenize{#168}}, \texttt{\detokenize{#170}}, \texttt{\detokenize{#171}}, \texttt{\detokenize{#172}}, \texttt{\detokenize{#173}}, \texttt{\detokenize{#174}}, \texttt{\detokenize{#177}}, \texttt{\detokenize{#208}} \\
M2 & S5-S8 & 164 & 37 & 14 & \texttt{\detokenize{#233}} \\
M3 & S9-S10 & 82 & 38 & 14 & Sin PR con milestone asignado en GitHub \\
M4 & S11-S13 & 277 & 132 & 20 & \texttt{\detokenize{#232}}, \texttt{\detokenize{#488}}, \texttt{\detokenize{#628}} \\
\end{longtable}
\endgroup
\end{landscape}

Los commits por milestone se calculan por ventana de sprint, no por enlace directo issue-commit. Cuando no hay PR con milestone asignado, se deja explícito para no inventar trazabilidad.

\section{8. Participación por integrante}

\subsection{Issues asignadas en M1-M4}

Las issues pueden tener más de una persona asignada; por eso la suma de asignaciones puede superar el total de issues.

\begin{landscape}
\begingroup
\scriptsize
\setlength{\LTpre}{0.4em}
\setlength{\LTpost}{0.8em}
\begin{longtable}{>{\RaggedRight\arraybackslash}p{0.153\linewidth}>{\RaggedRight\arraybackslash}p{0.153\linewidth}>{\RaggedRight\arraybackslash}p{0.153\linewidth}>{\RaggedRight\arraybackslash}p{0.153\linewidth}>{\RaggedRight\arraybackslash}p{0.153\linewidth}>{\RaggedRight\arraybackslash}p{0.153\linewidth}}
\toprule
\textbf{Integrante} & \textbf{M1 issues/HU} & \textbf{M2 issues/HU} & \textbf{M3 issues/HU} & \textbf{M4 issues/HU} & \textbf{Lectura} \\
\midrule
\endfirsthead
\toprule
\textbf{Integrante} & \textbf{M1 issues/HU} & \textbf{M2 issues/HU} & \textbf{M3 issues/HU} & \textbf{M4 issues/HU} & \textbf{Lectura} \\
\midrule
\endhead
\midrule
\multicolumn{6}{r}{\small\textit{Continua en la siguiente pagina}} \\
\endfoot
\bottomrule
\endlastfoot
\texttt{\detokenize{Sarm-m}} & 21 / 5 & 9 / 4 & 5 / 1 & 43 / 3 & Alta participación en coordinación, evidencia y cierre. \\
\texttt{\detokenize{samuelfl680}} & 16 / 2 & 3 / 2 & 11 / 5 & 56 / 10 & Alta carga en configuración, pruebas y cierre técnico. \\
\texttt{\detokenize{juanvargax}} & 42 / 12 & 5 / 1 & 10 / 2 & 22 / 2 & Alta participación en backlog, producto y planeación. \\
\texttt{\detokenize{solonlosada2006}} & 11 / 1 & 18 / 6 & 14 / 3 & 30 / 4 & Participación sostenida en DevOps, CI/CD y flota. \\
\texttt{\detokenize{Juserora}} & 21 / 4 & 6 / 1 & 11 / 3 & 17 / 2 & Participación en QA, pruebas, hallazgos y validación. \\
\end{longtable}
\endgroup
\end{landscape}

\subsection{Commits por integrante}

\begingroup
\small
\setlength{\LTpre}{0.4em}
\setlength{\LTpost}{0.8em}
\begin{longtable}{>{\RaggedRight\arraybackslash}p{0.313\linewidth}>{\RaggedRight\arraybackslash}p{0.313\linewidth}>{\RaggedRight\arraybackslash}p{0.313\linewidth}}
\toprule
\textbf{Integrante normalizado} & \textbf{Commits en corte académico} & \textbf{Lectura} \\
\midrule
\endfirsthead
\toprule
\textbf{Integrante normalizado} & \textbf{Commits en corte académico} & \textbf{Lectura} \\
\midrule
\endhead
\midrule
\multicolumn{3}{r}{\small\textit{Continua en la siguiente pagina}} \\
\endfoot
\bottomrule
\endlastfoot
\texttt{\detokenize{Sarm-m}} & 233 & Mayor volumen de commits, especialmente documentación, cierre y coordinación técnica. \\
\texttt{\detokenize{samuelfl680}} & 136 & Participación fuerte en configuración, pruebas y soporte técnico. \\
\texttt{\detokenize{juanvargax / juansebastianvd}} & 114 & Participación funcional y de producto con alias normalizados. \\
\texttt{\detokenize{solonlosada2006}} & 72 & Aportes asociados a DevOps, pipelines y despliegue. \\
\texttt{\detokenize{juserora}} & 51 & Aportes de QA, pruebas y ajustes funcionales. \\
Otros autores/bots & 24 & Cuentas de laboratorio, bot y autores sin normalización completa. \\
\end{longtable}
\endgroup

\section{9. Pull Requests y revisiones}

\begingroup
\small
\setlength{\LTpre}{0.4em}
\setlength{\LTpost}{0.8em}
\begin{longtable}{>{\RaggedRight\arraybackslash}p{0.313\linewidth}>{\RaggedRight\arraybackslash}p{0.313\linewidth}>{\RaggedRight\arraybackslash}p{0.313\linewidth}}
\toprule
\textbf{Métrica PR} & \textbf{Valor} & \textbf{Observación} \\
\midrule
\endfirsthead
\toprule
\textbf{Métrica PR} & \textbf{Valor} & \textbf{Observación} \\
\midrule
\endhead
\midrule
\multicolumn{3}{r}{\small\textit{Continua en la siguiente pagina}} \\
\endfoot
\bottomrule
\endlastfoot
PRs totales & 161 & Incluye todo el repositorio. \\
PRs mergeadas & 134 & Evidencia de integración por ramas. \\
PRs abiertas & 1 & \texttt{\detokenize{#628}}, M4, pendiente de integración. \\
PRs cerradas sin merge & 26 & Trabajo no integrado o reemplazado. \\
Reviews detectadas & 8 & Dato disponible desde \texttt{\detokenize{gh}}; no se inventan revisiones faltantes. \\
\end{longtable}
\endgroup

Reviews detectadas:

\begingroup
\small
\setlength{\LTpre}{0.4em}
\setlength{\LTpost}{0.8em}
\begin{longtable}{>{\RaggedRight\arraybackslash}p{0.47\linewidth}>{\RaggedRight\arraybackslash}p{0.47\linewidth}}
\toprule
\textbf{Reviewer} & \textbf{Reviews} \\
\midrule
\endfirsthead
\toprule
\textbf{Reviewer} & \textbf{Reviews} \\
\midrule
\endhead
\midrule
\multicolumn{2}{r}{\small\textit{Continua en la siguiente pagina}} \\
\endfoot
\bottomrule
\endlastfoot
\texttt{\detokenize{juanvargax}} & 3 \\
\texttt{\detokenize{samuelfl680}} & 2 \\
\texttt{\detokenize{copilot-pull-request-reviewer}} & 2 \\
\texttt{\detokenize{copilot-swe-agent}} & 1 \\
\end{longtable}
\endgroup

\section{10. Control de versiones, ramas y releases}

Ramas verificadas:

\texttt{\detokenize{develop}}, \texttt{\detokenize{feature-Slosada}}, \texttt{\detokenize{feature-sarm-m}}, \texttt{\detokenize{main}}, \texttt{\detokenize{origin/copilot/check-web-functionality-errors}}, \texttt{\detokenize{origin/develop}}, \texttt{\detokenize{origin/feature-Slosada}}, \texttt{\detokenize{origin/feature-Vargas-J}}, \texttt{\detokenize{origin/feature-juserora}}, \texttt{\detokenize{origin/feature-samuelfl680}}, \texttt{\detokenize{origin/feature-sarm-m}}, \texttt{\detokenize{origin/feature/docker-deploy}}, \texttt{\detokenize{origin/feature/pipelines}}, \texttt{\detokenize{origin/main}}, \texttt{\detokenize{origin/revert-422-feature-Slosada}}.

No existen tags/releases verificables. La trazabilidad por versión se sustenta con ramas, PRs, merges, capturas de GitHub, commits y milestones cerrados.

\section{11. Postmortem general del semestre}

\subsection{Producto}

Se produjo una plataforma MERN para gestión de flota, autenticación, documentos, alertas, métricas de calidad, SMS/Twilio, Docker Compose y CI/CD. El producto evolucionó desde una base funcional inicial hacia un sistema desplegable con evidencia de pruebas, SonarCloud, DockerHub y GitHub Actions.

\subsection{Esfuerzo}

El esfuerzo fue medible por issues, HU, commits, PRs y roles. La participación no fue idéntica por volumen, pero sí tuvo contribuciones claras por especialidad: Scrum, producto, configuración, DevOps y QA.

\subsection{Proceso}

El proceso mejoró de forma incremental: GitFlow y autonomía en M1, issues atómicas en M2, QA y Scrum Poker en M3, y cierre con CI/CD, SonarCloud y evidencia en M4.

\subsection{Roles}

La guía del profesor pide evaluar roles con hechos objetivos. Por eso los documentos de milestone conectan cada rol con evidencia: issues asignadas, commits, PRs, hallazgos de QA, configuración, CI/CD y seguimiento de riesgos.

\section{12. Causas raíz principales}

\begingroup
\scriptsize
\setlength{\LTpre}{0.4em}
\setlength{\LTpost}{0.8em}
\begin{longtable}{>{\RaggedRight\arraybackslash}p{0.235\linewidth}>{\RaggedRight\arraybackslash}p{0.235\linewidth}>{\RaggedRight\arraybackslash}p{0.235\linewidth}>{\RaggedRight\arraybackslash}p{0.235\linewidth}}
\toprule
\textbf{Problema} & \textbf{Causa raíz técnica} & \textbf{Causa raíz de proceso} & \textbf{Acción de mejora} \\
\midrule
\endfirsthead
\toprule
\textbf{Problema} & \textbf{Causa raíz técnica} & \textbf{Causa raíz de proceso} & \textbf{Acción de mejora} \\
\midrule
\endhead
\midrule
\multicolumn{4}{r}{\small\textit{Continua en la siguiente pagina}} \\
\endfoot
\bottomrule
\endlastfoot
Integraciones con fallos iniciales & Falta de validación local antes de merge & Definition of Done incompleto & Exigir build local y checklist antes de PR. \\
Retrabajo por entornos distintos & Rutas, alias e infraestructura no homogéneos & Riesgo técnico no identificado al planear & Usar Docker Compose y reglas de naming desde el inicio. \\
Credenciales expuestas temporalmente & Variables \texttt{\detokenize{.env}} configuradas tarde & Gestión de secretos tardía & Definir \texttt{\detokenize{.env.example}}, secrets y revisión de seguridad desde día 1. \\
Quality Gate bajo presión & Cobertura y hotspots tratados tarde & Calidad vista como cierre, no como práctica continua & Incrementar pruebas durante cada sprint. \\
Evidencia incompleta o tardía & Capturas y anexos no recolectados al cerrar HU & Evidencia no estaba en el DoD & Incorporar evidencia requerida antes de cerrar issues. \\
\end{longtable}
\endgroup

\section{13. Retrospectiva estrella de mar general}

\subsection{Comenzar a hacer}

\begin{itemize}[leftmargin=1.5em]
\item Definir un Definition of Done con evidencia, build local, revisión y captura mínima antes de cerrar cada HU.
\item Registrar datos de métricas ágiles al cierre de cada sprint para evitar reconciliaciones tardías.
\item Automatizar exportaciones con \texttt{\detokenize{gh}} y anexarlas al repositorio después de cada milestone.
\end{itemize}

\subsection{Más de}

\begin{itemize}[leftmargin=1.5em]
\item Revisar PRs con foco funcional y no solo con foco de compilación.
\item Usar Scrum Poker cuando haya historias con incertidumbre técnica.
\item Capturar evidencia visual en el mismo momento en que se valida una funcionalidad.
\end{itemize}

\subsection{Seguir haciendo}

\begin{itemize}[leftmargin=1.5em]
\item Mantener issues, PRs y commits trazables para sustentar trabajo individual y de equipo.
\item Conservar postmortems por milestone con causas raíz y acciones concretas.
\item Usar GitHub como fuente objetiva de seguimiento académico.
\end{itemize}

\subsection{Menos de}

\begin{itemize}[leftmargin=1.5em]
\item Depender de cierres masivos al final del sprint.
\item Usar métricas manuales sin explicar su criterio de cálculo.
\item Abrir PRs de cierre sin validación funcional local previa.
\end{itemize}

\subsection{Dejar de hacer}

\begin{itemize}[leftmargin=1.5em]
\item Contar milestones de taller como evidencia funcional del producto.
\item Presentar cifras sin separar issues, HU, PRs y milestone items.
\item Postergar seguridad, QA y evidencia para los últimos días.
\end{itemize}

\section{14. Relación explícita con la rúbrica}

\begingroup
\small
\setlength{\LTpre}{0.4em}
\setlength{\LTpost}{0.8em}
\begin{longtable}{>{\RaggedRight\arraybackslash}p{0.47\linewidth}>{\RaggedRight\arraybackslash}p{0.47\linewidth}}
\toprule
\textbf{Criterio de la rúbrica Agile/Postmortem} & \textbf{Evidencia en este paquete} \\
\midrule
\endfirsthead
\toprule
\textbf{Criterio de la rúbrica Agile/Postmortem} & \textbf{Evidencia en este paquete} \\
\midrule
\endhead
\midrule
\multicolumn{2}{r}{\small\textit{Continua en la siguiente pagina}} \\
\endfoot
\bottomrule
\endlastfoot
Milestones finalizados & M1-M4 documentados y cerrados en issues. \\
HU terminadas y trazables & 72 HU estrictas cerradas, con desglose por sprint y milestone. \\
HU por sprint y promedio & Tabla de 13 sprints y promedio 5.54 HU/sprint. \\
Commits por sprint y promedio & Tabla de 13 sprints y promedio 48.46 commits/sprint. \\
Distribución de participación & Issues/HU por integrante y commits normalizados. \\
PRs y revisiones & 161 PRs, 134 mergeadas, 8 reviews detectadas. \\
Control de versiones & Ramas, PRs, merges y ausencia explícita de releases/tags. \\
Trazabilidad por versión & PRs, ramas y evidencias de cierre; releases no verificables. \\
Postmortem profundo & Producto, esfuerzo, proceso, roles, causas raíz y acciones. \\
Acciones concretas & Planes de mejora por milestone y plan general. \\
\end{longtable}
\endgroup

\section{15. Métricas pendientes}

\begingroup
\small
\setlength{\LTpre}{0.4em}
\setlength{\LTpost}{0.8em}
\begin{longtable}{>{\RaggedRight\arraybackslash}p{0.47\linewidth}>{\RaggedRight\arraybackslash}p{0.47\linewidth}}
\toprule
\textbf{Dato} & \textbf{Estado} \\
\midrule
\endfirsthead
\toprule
\textbf{Dato} & \textbf{Estado} \\
\midrule
\endhead
\midrule
\multicolumn{2}{r}{\small\textit{Continua en la siguiente pagina}} \\
\endfoot
\bottomrule
\endlastfoot
Reconciliación final entre 88 HU históricas y 72 HU estrictas & Pendiente de auditoría manual del criterio original. \\
Reconciliación entre 375 issues históricas y 283 issues cerradas M1-M4 & Pendiente de auditoría de issues sin milestone o no funcionales. \\
Reviews completas por PR vía GraphQL & Pendiente de exportación desde GitHub API si se requiere mayor detalle. \\
Releases/tags & No existen releases/tags verificables. \\
\end{longtable}
\endgroup

\section{16. Conclusión}

El semestre evidencia una aplicación rigurosa y cuantitativa de Scrum académico. Los documentos nuevos en \texttt{\detokenize{docs/Agile/}} separan los datos verificables de los pendientes, eliminan milestones de taller, fortalecen el postmortem con la guía del profesor y dejan trazabilidad defendible para la rúbrica final.
\end{document}
